\chapter{Il runtime di funzione}
\label{cap:17-runtime-funzione}

I capitoli precedenti hanno descritto il control plane. Questo attraversa il
confine e descrive che cosa gira dentro il container di una funzione, quale
contratto lo lega al control plane, e perch\'e quel contratto ha la forma che ha.

\section{Anatomia del container}

Un container gestito da \nf{} contiene due strati.

\begin{center}
\small
\begin{adjustbox}{max width=\textwidth}
\begin{tabular}{@{}lp{9cm}@{}}
\toprule
\textbf{Strato} & \textbf{Ruolo} \\
\midrule
\code{watchdog} & Entrypoint del container. Binario Rust statico, circa 2\,MB,
distribuito su un'immagine \code{FROM scratch}. Possiede il ciclo di vita del
processo utente. \\
\emph{runtime} & Il processo dell'utente: un'applicazione Spring Boot, un
servizio FastAPI, un binario Go, uno script, qualunque cosa risponda a uno dei
tre protocolli. \\
\bottomrule
\end{tabular}
\end{adjustbox}
\end{center}

La separazione risolve un problema pratico. Il control plane vuole parlare un
solo protocollo HTTP con ogni funzione; gli autori di funzioni vogliono scrivere
in qualunque linguaggio, incluso uno che non ha un server HTTP. Il watchdog
adatta il secondo al primo.

\section{I tre protocolli di runtime}

L'enumerazione \code{RuntimeMode} (\cref{cap:05-contratti-dati}) dichiara come il
watchdog deve parlare con il processo figlio.

\begin{center}
\small
\begin{adjustbox}{max width=\textwidth}
\begin{tabular}{@{}lp{8.8cm}@{}}
\toprule
\textbf{Modo} & \textbf{Protocollo} \\
\midrule
\code{HTTP} & Il watchdog sonda \code{GET /health} finch\'e il runtime \`e
pronto, poi fa da reverse proxy su \code{POST /invoke} verso
\code{127.0.0.1:8081}. \\
\code{STDIO} & Il watchdog scrive il payload JSON su stdin del processo, legge
la risposta JSON da stdout, e ne raccoglie l'uscita. \\
\code{FILE} & Il watchdog scrive \code{/tmp/input.json}, avvia il processo con
le variabili \code{INPUT\_FILE} e \code{OUTPUT\_FILE}, e legge
\code{/tmp/output.json}. \\
\bottomrule
\end{tabular}
\end{adjustbox}
\end{center}

I tre modi coprono uno spettro di invasivit\`a decrescente. \code{HTTP} richiede
un server nel processo utente ed \`e il pi\`u efficiente, perch\'e il processo
resta vivo fra le invocazioni. \code{STDIO} non richiede nulla oltre alla
capacit\`a di leggere e scrivere JSON, ma paga la creazione di un processo per
invocazione. \code{FILE} \`e il ripiego per binari che non hanno un'interfaccia
migliore.

Il punto importante \`e che i tre modi sono \emph{indistinguibili dal control
plane}: qualunque sia il modo, il control plane invia una \code{POST} al
watchdog e riceve una risposta. La differenza \`e interamente confinata dietro
il watchdog.

\section{Warm e one-shot}

Il watchdog opera in due regimi.

In modalit\`a \code{DEPLOYMENT} resta vivo fra le invocazioni (container caldo):
espone il proprio server HTTP su \code{:8080}, esegue \code{WATCHDOG\_CMD} per
invocazione nei modi \code{STDIO} e \code{FILE}, e fa da reverse proxy verso il
runtime interno nel modo \code{HTTP}. Il control plane invia un
\code{InvocationRequest} come corpo e legge la risposta direttamente --- nessuna
callback per invocazione.

Il regime one-shot, precedente, avvia il processo una volta per invocazione e ne
riporta l'esito con una callback verso \code{CALLBACK\_URL}. \`E il regime che
paga un'inizializzazione per ogni richiesta, e la sua esistenza \`e ci\`o che
rende \code{coldStart} e \code{initDurationMs} campi significativi in
\code{ExecutionStatus}.

\subsection{Le variabili d'ambiente del contratto}

\begin{center}
\small
\begin{adjustbox}{max width=\textwidth}
\begin{tabular}{@{}llp{5.8cm}@{}}
\toprule
\textbf{Variabile} & \textbf{Regime} & \textbf{Significato} \\
\midrule
\code{FUNCTION\_NAME} & warm & Nome logico. \\
\code{WARM} & warm & \code{true}. \\
\code{TIMEOUT\_MS} & entrambi & Budget per invocazione. \\
\code{EXECUTION\_MODE} & entrambi & \\
\code{WATCHDOG\_CMD} & entrambi & Comando del processo utente. \\
\code{RUNTIME\_URL} & warm & \code{http://127.0.0.1:8081/invoke}. \\
\code{EXECUTION\_ID} & one-shot & Identificatore; nel warm arriva come header
\code{X-Execution-Id}. \\
\code{CALLBACK\_URL} & one-shot & Endpoint di completamento. \\
\code{INVOCATION\_PAYLOAD} & one-shot & Il payload, nell'ambiente. \\
\code{TRACE\_ID} & one-shot & \\
\bottomrule
\end{tabular}
\end{adjustbox}
\end{center}

Sono le stesse variabili che il builder Kubernetes protegge dalla sovrascrittura
da parte dell'utente (\cref{cap:15-deployment-provider}).

Il passaggio dal regime one-shot al warm si legge nella tabella: nel warm
l'identificatore di esecuzione si sposta dall'ambiente agli header, il che \`e
obbligato --- un processo che serve molte invocazioni non pu\`o riceverne
l'identit\`a in una variabile d'ambiente fissata all'avvio.

\section{Il contratto della risposta}

\subsection{\code{HandlerResponse}}

Un handler pu\`o restituire un valore semplice, che diventa l'output, oppure una
busta:

\begin{lstlisting}[language=Java]
/**
 * Optional response envelope a handler may return instead of a plain value to control
 * the HTTP status code, a safe set of response headers, and whether {@code output} is
 * base64-encoded binary. Its mere presence as the handler's return value is itself the
 * signal that the response was function-decided, not a platform error.
 */
public record HandlerResponse(Object output, int statusCode, Map<String, String> headers, String encoding) { ... }
\end{lstlisting}

La frase decisiva \`e l'ultima: \emph{la sola presenza della busta \`e il segnale
che la risposta \`e stata decisa dalla funzione}. \`E ci\`o che consente al
control plane di distinguere un \code{404} applicativo --- ``la risorsa che hai
chiesto alla mia funzione non esiste'' --- da un \code{404} di piattaforma ---
``questa funzione non \`e registrata''. La distinzione viaggia sul filo come
header \code{X-NanoFaaS-Function-Status} ed \`e interpretata sia dal dispatcher
esterno (\cref{cap:06-nucleo}) sia dal gateway di offload
(\cref{cap:13-offload}).

Senza questo meccanismo, il tasso di errore misurato da una piattaforma FaaS \`e
ambiguo: mescola i fallimenti della piattaforma con le risposte non-2xx
deliberate delle funzioni. Il \cref{cap:21-osservabilita} riprende il punto.

\subsection{La lista bianca degli header}

Una funzione non pu\`o impostare header arbitrari sulla risposta:

\begin{lstlisting}[language=Java]
public static final Set<String> ALLOWED_RESPONSE_HEADERS = Set.of(
        "content-type", "location", "cache-control", "etag",
        "content-disposition", "content-language", "retry-after", "vary");
\end{lstlisting}

Otto header, tutti descrittivi del contenuto o della sua cacheabilit\`a. Sono
esclusi gli header hop-by-hop, quelli che governano la connessione, e quelli che
il control plane usa per il proprio protocollo. Un handler che potesse impostare
\code{X-NanoFaaS-Function-Status} potrebbe mentire sulla natura della propria
risposta; uno che potesse impostare \code{Transfer-Encoding} potrebbe corrompere
la connessione.

Il filtro contiene due dettagli che vale la pena isolare, perch\'e riguardano
classi di bug reali.

\paragraph{Collisioni per maiuscole.}
\begin{lstlisting}[language=Java]
/**
 * At most one entry survives per header name, compared case-insensitively: HTTP header names
 * are case-insensitive, so emitting both {@code Content-Type} and {@code content-type} would put
 * two colliding entries on the response. The first occurrence in iteration order wins and keeps
 * its original casing — that casing is part of the public {@code InvocationResponse.headers}
 * contract and callers read it.
 */
\end{lstlisting}

\paragraph{Il locale turco.}
\begin{lstlisting}[language=Java]
// Locale.ROOT, not the default locale: under a Turkish/Azerbaijani default,
// 'I' folds to dotless 'i' (U+0131) and the allow-list match would silently fail.
String lowerKey = entry.getKey().toLowerCase(Locale.ROOT);
\end{lstlisting}

In turco la minuscola di \code{I} \`e \code{\i} senza punto, e
\code{"Location".toLowerCase()} su una JVM con locale turco non produce
\code{"location"}. Il confronto con la lista bianca fallirebbe, e l'header
sarebbe silenziosamente scartato --- su alcune macchine e non su altre. Lo stesso
accorgimento compare nel controller per il valore dell'header di ragione del
rifiuto (\cref{cap:04-architettura-control-plane}).

Sono dettagli, ma sono i dettagli che distinguono un contratto che funziona da
uno che funziona sulla macchina dello sviluppatore.

\section{Il watchdog}

Il watchdog \`e scritto in Rust, compilato staticamente con musl, e distribuito
su un'immagine \code{FROM scratch} di circa 2\,MB. Il sorgente principale \`e di
circa 1\,500 righe.

La scelta di Rust e della compilazione statica \`e coerente con il ruolo. Il
watchdog \`e presente in \emph{ogni} container di funzione, quindi la sua
impronta si moltiplica per il numero di repliche di ogni funzione; ed \`e
l'entrypoint, quindi il suo tempo di avvio si somma a ogni cold start. Un
binario statico da 2\,MB su un'immagine vuota \`e ci\`o che rende quei due costi
trascurabili.

Le sue responsabilit\`a: leggere la configurazione dall'ambiente, avviare il
processo figlio, attendere che diventi pronto, applicare il timeout per
invocazione, inoltrare la richiesta secondo il modo configurato, inviare la
callback nel regime one-shot, e gestire i segnali di terminazione.

\begin{damisurare}
Il costo del watchdog sul percorso di invocazione --- l'ulteriore salto di
reverse proxy nel modo \code{HTTP} --- non \`e stato isolato. \`E incluso in
tutte le misure di \code{function\_latency\_ms} riportate nel
\cref{cap:23-prestazioni-base}, ma non \`e stato confrontato con un'invocazione
diretta al runtime.
\todomis{latenza aggiunta dal reverse proxy del watchdog}
\end{damisurare}

\section{Il servizio di riferimento \code{warm-echo}}

\code{services/java/warm-echo} \`e un servizio che incorpora l'SDK Java e si
limita a restituire l'input. Sono ventiquattro righe.

La sua funzione nel progetto \`e sperimentale: fornisce un runtime \emph{reale},
con un vero salto HTTP e un vero processo JVM, ma con un tempo di esecuzione
dell'handler prossimo a zero. Il confronto con la modalit\`a \code{LOCAL}
(\cref{cap:07-modalita-esecuzione}) isola cos\`i il costo del trasporto dal costo
del control plane.

Il progetto insiste su un punto: \code{warm-echo} \`e distribuito attraverso la
normale API di registrazione, non installato come infrastruttura condivisa. La
distinzione conta perch\'e mantiene il servizio di riferimento soggetto alle
stesse politiche --- code, limiti di concorrenza, retry --- di qualunque altra
funzione. Un servizio di riferimento privilegiato misurerebbe un percorso che
nessun'altra funzione percorre.
